\documentclass[12pt]{article}

\usepackage[utf8]{inputenc}
\usepackage[T1]{fontenc}

\title{El lenguaje científico de la Biología}
\author{Kihara G. Rendon, Sara S. Pay, Diego A. Castro, Alessandro I. Molina}
\date{11/09/2026}

\begin{document}

\maketitle
\section{Autores del artículo de referencia}
\textbf{Artículo de referencia:} Bazán Delgado Anieska. (2017). \textit{LA BIOLOGÍA Y SU LENGUAJE CIENTÍFICO, UN ANÁLISIS LINGüÍSCO}. Revista Santiago.

\section{Introducción}
En la formación escolar y universitaria, el aprendizaje de las ciencias ha enfrentado barreras como la memorización pasiva o la falta de motivación. La autora aborda este problema centrándose en la complejidad del vocabulario científico. En su trabajo, integró los aspectos conceptuales y léxicos que diferenciaban este discurso del habla cotidiana del resto de disciplinas. Para quienes nos formamos en el área de la Biología, este enfoque resulta clave para comprender que el manejo preciso de la terminología no es un simple formalismo, sino es la herramienta fundamental para construir un pensamiento científico. Cuando se asume el lenguaje técnico como parte activa del razonamiento y no como una lista de palabras aisladas, la forma en que los estudiantes se enfrentan a los contenidos de la carrera cambia por completo. 

En disciplina donde los procesos son tan especificos y dinámicos como en las ciencias biológicas, entender la raiz y el contexto de cada término permite conectar fenómenos que, a primera vista, parecen separados. De este modo, el estudiante deja de repetir definiciones de memoria para empezar interpretar lo que ocurre realmente, logrando apropiarse del conocimiento desde una base mucho mas sólida y reflexiva.

Por otro lado, esta integración entre el léxico y el concepto tiene un impacto directo en la manera en que los futuros profesionales logran comunicar sus observaciones. A menudo, la dificultad para redactar un buen informe de laboratorio o discutir un resultado no radica en no haber entendido la práctica, sino en no contar con las herramientas lingüisticas para expresarlo con rigor. Al trabajar de manera consciente la estructura de las palabras y su uso adecuado dentro del discurso científico, se reduce la brecha entre lo que el alumno comprende mentalmente y lo que es capaz de plasmar por escrito, fortaleciendo asi su desempeño académico desde los primeros semestres.

\section{El lenguaje como herramienta para construir el conocimiento biológico}

El lenguaje científico cumple una función esencial en el desarrollo de la biología, debido a que permite que se exprese conocimiento de manera precisa y comprensible dentro de una comunidad académica. A diferencia del leguaje común, que puede presentar expresiones ambiguas o depender del contexto de una conversación, el lenguaje científico busca establecer significados específicos que faciliten la interpretación de los fenómenos estudiados. En biología, esta precisión  es observada en la necesidad de diferenciar conceptos que pueden parecer similares, pero que representan estructuras diferentes. por ejemplo, términos como mitosis, meiosis, difusión y ósmosis no son intercambiables, ya que cada uno describe un fenómeno particular. Comprender sus diferencias permite interpretar adecuadamente los procesos biológicos y establecer reaciones entre los contenidos de la carrera.

El lenguje tambien permite organizar ideas y convertirlas en explicaciones que puedan ser comprendidas y discutidas por sus investigadores. Cuando un biólogo describe una estructura celular, identifica una especie o explica una interacción ecológica, necesita utilizar conceptos que representen aquello que está estudiando. De forma que las palabras no funcionan simplemente como instrumentos que permiten clasificar, comparar y relacionar la información onbtenida mediante la experimentación y la capacidad de observar.

Es importante reconocer que el lenguaje en la biología no permanece quieto. A medida que surgen nuevos descubrimientos y se desarrollan campos de investigación, aparecen conceptos y términos que permiten explicar fenómenos antes desconocidos, Esto demuestra que el vocabulario científico se transforma junto con el conocimiento y que su aprendizaje requiere uan comprensión continua de los conceptos. Por esta razón, el dominio del lenguaje especializado no debe de estar limitado a la repeteción de definiciones. Un estudiante realmente comprende un término cuando  puede reconocer a que fenómeno hace referencia, relacionarlo con otros conocimientos y utilizarlo correctamente para explicar una situación biológica.

\section{Formación del vocabulario científico}

Una de las características ms representativas del lenguaje biológico es la presencia de términos construidos apartir de el latin y  el griego. Estos elemnetos han contribuido a la formación de gran parte del vocabulario utilizado para describir organismos, estructuras, funciones y procesos. Su conocimiento puede facilitar la comprensión de palabras desconocidas, debido a que permite reconocer los compmentes que participan en su formación. Por ejemplo, la fotosíntesis contiene elementos relacionados a la luz y la sintesis o formación de sustancias. De igual, manera la palabra morfología se relaciona con el estudio de las formas y las estructuras. Estos ejemplos muestran cómo el reconocimiento de palabras desconocidas permite aproximarse al significado de los términos sin depender exclusivamente de su memorización.

El analisis etimológico puede convertirse en una gran herramienta de aprendizaje.Ya que genera la posibilidad de establecer relaciones con otros términos que comparten la misma raíz. Asi conceptos como biología, biodiversidad y biosfera presentan un elemento común emparentado con la vida, este reconocimiento facilita la construcción de asociaciones y contribuyen a ampliar progresivamente el vocabulario científico. La biología incorpora términos provenientes de otras disciplinas como la química, la física y la geografía; bioquimica, biofísica, biogeografía, cuyos conceptos integran conocimientos de áreas diferentes para estudiar fenómenos relacionados con los seres vivos. El uso de abreviaciones y simbolos en diferentes áreas de la biología tales como: ARN, ADN, pH y ATP, tambien permiten comunicar información de manera breve y reconocible. Sin embargo, su correcta interpretación requiere concocer el significado y el contexto en el que se utilizan. Por consiguiente, comprender cómo se forma el vocabulario biologíco permite superar la idea de que los términos científicos son palabras aisladas que deben aprendersen de memoria. Su estructura contiene información que puede orientar la comprensión de conceptos y facilitar el aprendizaje.

\section{La nomenclatura y la precisión en la comunicación biológica}
La necesidad de identificar y clasificar los organismos ha llevado al desarrollo de sistemas de nomenclatura que permiten establecer formas de denominación reconocidas por la comunidad científica. Uno de los aportes más importantes en este campo fue el trabajo de Carlos Linneo, quien contribuyó a consilidar el sistema de nomenclatura binomial utilizado para nombrar las especies mediante el género y el epíteto específico. Este sistea permite que los organismos puedan identificarse de manera más precisa, independientemente de las diferencias lingüísticas. Por otra parte, en el lenguaje científico también se encuentra los epónimos, que corresponden a términos derivados de nombres propios. De esta manera, la nomenclatura constituye una herramienta que permite organizar la diversidad biológica, estableciendo criterios comunes para la identificación de los organismos.

\section{La estructura del lenguaje en los textos científicos}
El lenguaje científico no se caracteriza únicamente por el uso de términos especializados, sino también por una formar particular de organizar las ideas, donde se es frecuente encontrar contrucciones gramaticales que permiten presentar la información de manera objetiva, precisa y organizada. Entre estos recursos se encuentra el uso de construcciones impersonales y de partícula, que permiten centrar la atención en los procesos y resultados, en lugar de destacar a quien realiza la acción. 

Asimismo, las oraciones pueden incluir estructuras complejas que relacionan diferentes conceptos y permiten explicar procesos mediante relaciones de causa, consecuencia y comparación. También se emplean formas no personales del verbo, como el infinitivo, el gerundio y el participio, que contribuyen a organizar la información y expresar las acciones de manera diferente. Otro aspecto relevante es la nominalización, mediante la cual una acción o un proceso se expresa a través de un sustantivo. Este recurso permite condesar información y es frecuente en el discursi académico. Sin embargo, cuando se utiliza en exceso, puede dificultar la lectura, especialmente para quines todavía están desarrollando habilidades de comprensión.

\section{ La importancia del lenguaje científico en la formación de los biólogos}
El dominio del lenguaje tiene consecuencias directas en la formación académica y profesional de los estudiantes de biología. Durante la carrera, se es necesario interpretar artículos, describir organismos, registrar observaciones y comunicar resultados. Todas estas actividades requieren utilizar términos adecuados para contruir explicaciones que representen correctamente los fenómenos estudiados. En ocasiones, las dificultades para redactar o explicar un resultado no se deben únicamente a la falta de comprensión sino también a las limitaciones para expresar las ideas con claridad y precisión. Esto evidencia que la comprensión conceptual y la capacidad de comunicación se encuentran estrechamente relacionadas.

En este contexto, la enseñanza de la biología podría incorporar actividades orientadas al análisis del vocabulario y de las estructuras del discurso. De igual manera, la elaboración de descripciones y la redacción de pequeños textos científicos desde los primeros semestres permitirían fortalecer progresivamente las habilidades de comunicacón. Estas actividades podrían complementarse con la lectura de artículos científicos, en lo que se van identificando la terminología utilizada y se reconozcan cómo se construyen las explicaciones. 
No se trata de convertir la enseñanza en clases de lingüística ni de restarle importancia a otros aspectos de la carrera. Por el contrario, se busca integrar el lenguaje como una herramienta de comprensión y expresión de los fenómenos biológicos y asi el estudiante adquiera una mayor confianza para expresar sus ideas.

\section{Conclusión}
La lingüística de la biología revela una dimensión del conocimiento científico que con frecuencia queda oculta detrás de los contenidos. La biología no se aprende únicamente mediante la observación de organismos, la experimentación o el estudio de modelos; también se aprende al dominar un sistema de palabras, estructuras, símbolos, convenciones y formas de argumetar. Ese sistema permite nombrar entidades, clasificar formas de vida, representar procesos, formular relaciones causales, comparar evidencias, y construir explicaciones. Por ello, el lenguaje no es un envoltorio externo de la ciencia biológica: es uno de los medios a través de los cuales la ciencia adquiere forma y puede ser compartida.

El análisis del lenguaje muestra que la terminologìa biológica contiene patrones morfológicos, raíces grecolatinas, derivaciones, abreviaciones y préstamos interdisciplinarios que pueden convertirse en recursos pedagógicos. La evidencia de Brown (2014) y Keen et al. (2025) sugiere que enseñar a interpretar la estructura de las palabras puede aportar beneficios al aprendizaje, aunque no sustituye la enseñanza conceptual. Del mismo modo, los estudios de Rector et al. (2013) muestran que la ambigüedad léxica puede afectar las explicaciones de evolución, mientras que trabajos sobre sintaxis y nominalización evidencian que comprender la gramática del discurso científico es necesario para leer textos biológicos complejos (Hao & Humphrey, 2019; Halliday & Martin, 1993).

La principal consecuencia de este análisis es pedagógica. Memorizar términos sigue siendo útil e incluso necesario en una disciplina tan extensa como la biología, pero memorizar no equivale a comprender. La terminología adquiere verdadero valor cuando el estudiante puede reconocer su estructura, diferenciar sus usos, relacionarla con otros conceptos, utilizarla en una explicación y trasladarla entre diferentes representaciones. La alfabetización científica, por tanto, incluye una alfabetización lingüística disciplinar.

En este marco, los talleres de análisis etimológico, contraste entre lenguaje cotidiano y científico, desempaque de nominalizaciones, construción de textos y discusión de conceptos pueden integrarse de forma progresiva en los cursos de biología. No se trata de agregar una asignatura aislada que compita con los contenidos biológicos, sino de transformar la manera en que esos contenidos se enseñan. El lenguaje debe aparecer en la clase de biología como parte del fenómeno educativo mismo.

Finalmente, comprender la lingüística de la biología también significa comprender mejor la naturaleza de la ciencia. La terminología cambia porque cambian las preguntas científicas; las nomenclaturas se regulan porque las disciplinas aumulan conocimiento; y las palabras adquieren significados especializados porque cada campo construye sus propios modelos de realidad. Enseñar a los estudiantes a leer y producir ese lenguaje es, en consecuencia, enseñarles a participar de la cultura científica.
\end{document}

archivo con actividad
